% !PRINTED 13
\weberpage{13}
%<<CONT
von zwei Verhältnissen mit demselben Zähler und ungleichen Nennern ist
das das kleinere, dessen Nenner grösser ist.

Sind nämlich $a,a',b$ beliebige Elemente einer messbaren Menge und
$a'>a$, so wähle man zunächst eine ganze Zahl $n$ so, dass $na>b$ und
$n(a'-a)>b$. Hierauf nehme man die \gesperrt{kleinste} ganze Zahl $m$,
die der Bedingung $mb>na$ genügt; dann ist $na<mb$, aber $mb<na'$.
% [sic] Der Druck setzt in der folgenden Kette auch zwischen na' und
% na+n(a'-a) das Zeichen \ge, obwohl dort algebraisch Gleichheit gilt.
Denn wäre $mb\ge na'\ge na+n(a'-a)$, so wäre $mb>na+b$; also wäre gegen
die Voraussetzung schon $(m-1)b>na$. Dann ist
\[
 {a\over b}<{m\over n}<{a'\over b},
\]
also $a:b<a':b$; und ganz ebenso kann man beweisen, dass, wenn $b'>b$
ist, $a:b>a:b'$ wird.

Man drückt diesen Satz auch so aus, dass ein Verhältniss zugleich mit
dem Zähler wächst und mit wachsendem Nenner abnimmt.

Hieraus ergiebt sich auch leicht der folgende Satz: Sind $a,b,c,d$
Elemente derselben messbaren Menge und ist $a:b=c:d$, so ist auch
$a:c=b:d$.

Denn angenommen, es wäre $a:c<b:d$, so müsste es zwei ganze Zahlen
$m,n$ geben, so dass
\[
 na<mc,
 \qquad nb>md.
\]
Dann aber wäre nach dem eben bewiesenen Satze $na:nb<mc:md$, also auch
$a:b<c:d$, entgegen der Voraussetzung.

Hieran schliesst sich nun folgender \gesperrt{Hauptsatz}.
\gesperrt{Wenn von den vier Grössen} $a,b,c,d$ \gesperrt{irgend drei aus
einer stetigen messbaren Menge beliebig gegeben sind, so lässt sich das
vierte in derselben Menge so bestimmen, dass} $a:b=c:d$ \gesperrt{ist.}

Der Satz ist eine unmittelbare Folge der vorausgesetzten Stetigkeit.
Denn wenn man z.~B. ein Element $x$ von $\mathfrak M$ in $A$ oder in
$B$ aufnimmt, je nachdem $x:b$ kleiner oder grösser als $c:d$ ist, so
erhält man einen Schnitt, der durch ein Element $a$ erzeugt wird, das
der Bedingung $a:b=c:d$ genügt. Es gilt aber dieser Satz auch in
gewissen nicht stetigen Mengen, z.~B. für die rationalen Brüche.

Hieraus folgt, dass man als Repräsentanten zweier Zahlen immer zwei
Verhältnisse wählen kann, deren Elemente derselben
%CONT>>
